\section{Related work}
\label{sec:related}

\paragraph{Vision world models}
Following the taxonomy of Yu \etal~\cite{yu2026seeing}, vision world models fall into four design families: sequential
token generation~\cite{bruce2024genie}, diffusion-based generation~\cite{agarwal2025cosmos,bar2024nwm}, embedding
prediction~\cite{lecun2022path,bardes2024vjepa,assran2025vjepa2,zhou2024dinowm,baldassarre2025dinoworld}, and latent
state transition~\cite{ha2018world,hafner2023dreamerv3,micheli2023iris,zhang2023storm,robine2023twm}. Our student is an
embedding-prediction model with a state-transition core: a frozen \vjepa encoder supplies the target space, as DINO-WM
and DINO-world do with frozen DINO features~\cite{zhou2024dinowm,baldassarre2025dinoworld}, and a recurrent network
carries the latent through time under actions, as the recurrent state-space models of Dreamer~\cite{hafner2023dreamerv3}
and the transformer cores of IRIS, STORM and TWM~\cite{micheli2023iris,zhang2023storm,robine2023twm} do. The field's
recurrent core is a free architectural choice sized by compute; ours is fixed by a wiring diagram. We therefore compare
against the two cores a reviewer would ask for by name, a gated recurrent unit and a causal transformer predictor, both
matched to the student's parameter budget, and against the trivial copy-last and constant predictors every world-model
benchmark reports~\cite{bar2024nwm}. Knowledge distillation~\cite{hinton2015distilling} usually compresses a model into
a smaller member of the same family; here the target is a fixed, signed, biologically sourced sparse graph.

\paragraph{Navigation world models and embodied evaluation}
Navigation World Models~\cite{bar2024nwm} plan trajectories in a diffusion video model and evaluate with the
success-rate and SPL protocol of embodied navigation~\cite{anderson2018evaluation}; DINO-WM plans zero-shot in a frozen
latent space~\cite{zhou2024dinowm}. We adopt the same metrics for blind homing and goal-vector navigation, and the same
distinction between open-loop prediction and closed-loop control, which is where our behavioural results stop short.

\paragraph{Connectome-constrained models}
Whole-brain \emph{Drosophila} connectomes~\cite{dorkenwald2024neuronal,schlegel2024whole} with predicted neurotransmitter
signs~\cite{eckstein2024neurotransmitter} have enabled models whose connectivity is read off the wiring diagram rather than
learned. Lappalainen \etal~\cite{lappalainen2024connectome} fit a connectome-constrained visual system to a motion task
and predicted neural recordings; models of the central complex~\cite{hulse2021connectome,stone2017anatomically}
explain heading and path-integration circuits mechanistically. These works fit neural data or hand-designed tasks. We
instead distil a self-supervised video model into the wiring and measure what the wiring does to a world model,
with matched-control analysis as the main instrument.

\paragraph{Path integration in artificial networks}
Recurrent networks trained to localise from velocity input develop grid- and head-direction-like
codes~\cite{cueva2018emergence,banino2018vector,sorscher2019unified}. In those studies the task \emph{is} localisation.
Here the task is latent video prediction; heading and position are never supervised, and the probe results show that
whether such a code forms depends on the recurrent topology, not only on the objective.

\paragraph{Controls for wiring identity}
Degree-preserving rewiring~\cite{maslov2002specificity} is the standard null model for network structure. We extend it
with a block-preserving variant that keeps within-region degree sequences and the region-to-region edge counts, so that
the comparison isolates synapse-level identity from the macro-architecture, and with sign controls that keep Dale's law
while randomising which neuron is excitatory.
